\section{Chaîne probatoire continue de la publication exceptionnelle : du déploiement APP à Paley--Witt,
\texorpdfstring{\(V^\natural\)}{V natural} et Moonshine}
\label{sec:prueba-corredor-excepcional}

La publication exceptionnelle part du lecteur d'incidence
\(\mathfrak E_{\rm exc}\) de \eqref{eq:frontal-doble-proyeccion}. Elle ne prend pas en
entrée un code, un réseau ou un groupe exceptionnel déjà donnés : elle prend le
bord, l'orientation, le feuillet et la mémoire d'une histoire enrichie et
les exprime dans le calibre projectif fixé par l'holonomie nonadique. Ce n'est
qu'ensuite que les théorèmes classiques des codes, des réseaux et des algèbres
d'opérateurs de vertex sont appliqués aux objets produits.

\subsection{\texorpdfstring{Déploiement APP et cône gradué de l'indice \(54\)}{Déploiement APP et cône gradué de l'indice 54}}

Soit \(R\) l'involution qui échange les feuillets APP et soient
\(P_+=(I+R)/2\), \(P_-=(I-R)/2\) ses projecteurs. Le bloc central
\begin{equation}
 B_c=7I+2R=9P_++5P_-
 \label{eq:exc-bloque-central-app}
\end{equation}
produit quatre objets de types distincts :
\begin{equation}
 9-5=4,\qquad (B_c)_{i\ne j}=2,\qquad
 \sum M_\Pi-\sum M_\Sigma=51-45=6,\qquad
 9\cdot6=459-405=54.
 \label{eq:exc-despliegue-42654}
\end{equation}
Le \(4\) est le saut spectral entre feuillets, le \(2\) le couplage local,
le \(6\) la différence des compressions modulaires et le \(54\)
l'intégration de cette différence sur les neuf positions. Par conséquent,
\(4\to2\to6\to54\) désigne une hiérarchie de morphismes et d'évaluations, non
une chaîne d'égalités entre objets.

L'orientation \(o_{\Sigma\Pi}\), la parité de feuillet et la
normalisation triangulaire étant fixées, la branche de Fock possède le couplage
\begin{equation}
 \mathcal H_{\Sigma,\Pi}^{(m)}:
 M_{\rm bal}^{C_3,-}\longrightarrow
 \operatorname{End}\!\bigl(\mathcal F_2(\mathcal V_m)\bigr)
 \otimes o_{\Sigma\Pi},
 \label{eq:exc-acoplamiento-fock}
\end{equation}
et sa trace orientée satisfait
\begin{equation}
 \operatorname{Tr}_{o,\mathcal F_2}
 \bigl(\mathcal H_{\Sigma,\Pi}^{(m)}(\nu)\Gamma_2(g_m)\bigr)
 =-\frac{3m\,c(\nu)}2.
 \label{eq:exc-traza-fock}
\end{equation}
Pour \(m=12\) et \(c(\delta)=-3\), on a \(54\). Indépendamment, le
produit cyclotomique
\begin{equation}
 t_3(q)=q^{-1}\prod_{n\ge1}(1+q^n+q^{2n})^{-12}
 =q^{-1}-12+54q-76q^2-243q^3+\cdots
 \label{eq:exc-t3}
\end{equation}
donne \([q]t_3=54\).

\begin{theorem}[Cône typé de l'indice commun]
\label{thm:exc-cono-54}
Une fois construits séparément le recensement APP, la branche de Fock, le produit
cyclotomique, le voisin réticulaire de
\cref{eq:exc-vecino-indice-tres} et le module lunaire de
\cref{eq:exc-flm}, leurs lecteurs vérifient
\begin{equation}
 D_{\rm APP}
 =\operatorname{Tr}_{o,\mathcal F_2}
  \bigl(\mathcal H_{\Sigma,\Pi}^{(12)}(\delta)\Gamma_2(g_{12})\bigr)
 =[q]t_3=v^2=\frac{196884}{5\cdot729+1}=54.
 \label{eq:exc-cono-54}
\end{equation}
L'égalité réunit des évaluations dans un cône vers \(\mathbb Z\) ; elle
n'identifie pas leurs domaines et n'utilise pas \(196884\) pour sélectionner le corridor.
\end{theorem}

\begin{proof}
Les trois premières égalités sont
\eqref{eq:exc-despliegue-42654}, \eqref{eq:exc-traza-fock} et
\eqref{eq:exc-t3}. Pour le vecteur radial
\(v=4\rho_1+\rho_2+\cdots+\rho_{12}\) de
\eqref{eq:exc-vector-vecino}, l'orthogonalité des facteurs \(A_2\) donne
\(v^2=4^2\cdot2+11\cdot2=54\). Enfin,
\(5\cdot729+1=3646\) et \(54\cdot3646=196884\). Chaque identité est prouvée dans
son propre codomaine avant leur comparaison.
\end{proof}

\subsection{\texorpdfstring{Forme de Paley et racine quadratique interne de \(-I_6\)}{Forme de Paley et racine quadratique interne de -I6}}

Soit
\[
 \mathcal U_9=(\mathbb Z/9\mathbb Z)^\times=\{1,2,4,8,7,5\},
\]
dans l'ordre des puissances de \(2\), et soit
\begin{equation}
 \theta:\mathcal U_9\xrightarrow{\sim}\mathbb P^1(\mathbb F_5),
 \qquad
 \begin{array}{c|cccccc}
 u&1&2&4&8&7&5\\ \hline
 \theta(u)&\infty&0&1&2&3&4
 \end{array}.
 \label{eq:exc-calibre-proyectivo}
\end{equation}
Cette bijection est un calibre de coordonnées du lecteur ; elle n'est pas utilisée comme
homomorphisme. Pour le caractère quadratique
\(\chi:\mathbb F_5\to\{-1,0,1\}\), définissons
\begin{equation}
 C_P=
 \begin{pmatrix}
 0& 1& 1& 1& 1& 1\\
 1& 0& 1&-1&-1& 1\\
 1& 1& 0& 1&-1&-1\\
 1&-1& 1& 0& 1&-1\\
 1&-1&-1& 1& 0& 1\\
 1& 1&-1&-1& 1& 0
 \end{pmatrix},
 \label{eq:exc-matriz-paley}
\end{equation}
où les entrées finies hors diagonale sont \(\chi(x-y)\) et les
entrées incidentes à \(\infty\) valent un.

\begin{lemma}[Corrélation quadratique]
\label{lem:exc-correlacion-cuadratica}
Si \(a\ne b\) dans \(\mathbb F_5\), alors
\[
 \sum_{t\in\mathbb F_5}\chi(t-a)\chi(t-b)=-1.
\]
\end{lemma}

\begin{proof}
La substitution \(u=(t-a)/(b-a)\) réduit la somme à
\(\sum_u\chi(u(u-1))\). Pour \(u=0,1,2,3,4\), les termes sont
\(0,0,1,-1,-1\).
\end{proof}

\begin{proposition}[Identité de conférence et réduction ternaire]
\label{prop:exc-paley-witt}
La matrice \(C_P\) est symétrique et satisfait
\[
 C_PC_P^{\mathsf T}=5I_6.
\]
Sa réduction \(A_W=C_P\bmod3\) vérifie
\begin{equation}
 \boxed{A_W^2=-I_6\quad\text{dans }M_6(\mathbb F_3).}
 \label{eq:exc-aw-cuadrado}
\end{equation}
\end{proposition}

\begin{proof}
Chaque ligne de \(C_P\) a une norme au carré égale à cinq. Deux lignes finies
distinctes ont pour produit scalaire \(1-1=0\) par le
\cref{lem:exc-correlacion-cuadratica} ; la somme d'un caractère non trivial
prouve aussi l'orthogonalité avec la ligne de \(\infty\). En réduisant modulo
trois, on obtient
\(A_WA_W^{\mathsf T}=5I_6=2I_6=-I_6\). Comme \(A_W\) est symétrique,
\eqref{eq:exc-aw-cuadrado} en découle.
\end{proof}

L'identité \eqref{eq:exc-aw-cuadrado} construit la structure complexe
finie du corridor. En effet,
\[
 \mathbb F_9=\mathbb F_3[\iota]/(\iota^2+1),
 \qquad
 P_{\pm\iota}=\frac12(I\mp\iota A_W)
\]
sont des projecteurs complémentaires de rang trois et munissent
\(\mathbb F_3^6\) d'une structure de \(\mathbb F_9\)-module de rang trois.
L'action est explicite :
\begin{equation}
 (a+b\iota)\cdot v=av+b(vA_W).
 \label{eq:exc-accion-f9}
\end{equation}
Par conséquent, \(3^6=9^3\) n'est pas utilisé comme substitut cardinal d'une action.

\subsection{Code ternaire autodual et système de Witt}

Définissons
\begin{equation}
 c:\mathbb F_3^6\longrightarrow\mathbb F_3^{12},
 \qquad c(q)=(q,qA_W),
 \qquad C_W:=\operatorname{im}c.
 \label{eq:exc-codigo-grafico}
\end{equation}

\begin{theorem}[Code ternaire produit par l'incidence]
\label{thm:exc-codigo-ternario}
L'application \(c\) est injective et
\begin{equation}
 C_W=C_W^\perp,
 \qquad
 W_{C_W}(y)=1+264y^6+440y^9+24y^{12}.
 \label{eq:exc-enumerador-cw}
\end{equation}
En particulier, \(C_W\) a pour paramètres \([12,6,6]_3\).
\end{theorem}

\begin{proof}
La première projection de \(c(q)\) est \(q\), donc \(c\) est injective.
Avec \(G_W=(I_6\mid A_W)\),
\[
 G_WG_W^{\mathsf T}=I_6+A_WA_W^{\mathsf T}=0
\]
dans \(\mathbb F_3\). Ainsi \(C_W\subseteq C_W^\perp\), et l'égalité des
dimensions donne l'autodualité. L'énumérateur est la somme finie exacte
\[
 \sum_{q\in\mathbb F_3^6}
 y^{\operatorname{wt}(q)+\operatorname{wt}(qA_W)}.
\]
Évaluer les \(3^6=729\) lignes de la matrice explicite
\eqref{eq:exc-matriz-paley} donne, respectivement, les effectifs
\(1,264,440,24\) aux poids \(0,6,9,12\), et zéro aux autres. La somme
des quatre effectifs est \(729\), de sorte qu'aucun mot ne reste hors
du certificat fini.
\end{proof}

Les \(264\) vecteurs de poids six apparaissent par paires \(\{x,-x\}\) ; leurs
\(132\) supports forment le système de Steiner
\begin{equation}
 W_{12}=S(5,6,12).
 \label{eq:exc-w12}
\end{equation}
C'est la reconnaissance, après construction de \(C_W\), du code ternaire
étendu de Golay et de son plan de Witt
\citep{MacWilliamsSloane1977,ConwaySloane1999}. Le plan n'est pas introduit
pour sélectionner la matrice \(A_W\).

Le prolongement binaire demeure séparé. Si
\(\mathcal G_{24}\subset\mathbb F_2^{24}\) est le code binaire étendu
de Golay, \(D\) une dodécade et
\(\ell:\{1,\ldots,12\}_{\rm HMT}\to D\) un isomorphisme de plans, les
supports de poids huit forment \(W_{24}=S(5,8,24)\). Pour une tétrade fixe,
\[
 \lambda_4(W_{12})=4,
 \qquad
 \lambda_4(W_{24})=5,
\]
de sorte que quatre octades relèvent les quatre hexades incidentes et qu'une cinquième
est transversale. La donnée \((\mathcal G_{24},D,\ell)\) appartient au codomaine
binaire ; il n'existe pas de flèche canonique \(W_{24}\to\Lambda_{24}\).

\subsection{Recollement ternaire et voisin de Leech}

Soit \(Q=A_2^{12}\), avec \(\det Q=3^{12}\). Le discriminant de chaque
facteur \(A_2\) est \(\mathbb F_3\), et le recollement défini par
\eqref{eq:exc-codigo-grafico} est
\begin{equation}
 N:=\{x\in Q^*:x\bmod Q\in C_W\}.
 \label{eq:exc-niemeier-pegado}
\end{equation}

\begin{theorem}[Réseau de Niemeier produit par le code]
\label{thm:exc-niemeier}
Le réseau \(N\) est défini positif, pair et unimodulaire de rang
vingt-quatre, et son système de racines est \(A_2^{12}\). Par conséquent,
\(N\cong N(A_2^{12})\).
\end{theorem}

\begin{proof}
L'autodualité de \(C_W\) rend le recollement entier. Tous ses poids sont
des multiples de trois, ce qui rend l'extension paire. En outre,
\([N:Q]=|C_W|=3^6\), et donc
\[
 \det N=\frac{\det Q}{[N:Q]^2}
 =\frac{3^{12}}{3^{12}}=1.
\]
Une classe discriminante non nulle de \(A_2\) a pour norme minimale \(2/3\).
Comme tout mot non nul de \(C_W\) a un poids au moins égal à six, tout vecteur
de recollement extérieur à \(Q\) a une norme au moins égale à quatre. Les vecteurs de
norme deux sont exactement les racines des douze facteurs \(A_2\). La
classification des réseaux pairs unimodulaires de rang vingt-quatre
identifie le type indiqué \citep{ConwaySloane1999}.
\end{proof}

Soit \(\rho_i\) la plus haute racine du facteur \(A_2^{(i)}\), normalisée par
\(\rho_i^2=2\), et définissons
\begin{equation}
 v=4\rho_1+\rho_2+\cdots+\rho_{12},
 \qquad v^2=54.
 \label{eq:exc-vector-vecino}
\end{equation}
Le repère ordonné du lecteur exceptionnel fixe laquelle des douze fibres occupe
la première position ; aucune valeur archimédienne n'intervient dans cette sélection.
Posons
\begin{equation}
 \chi_v(x)=\langle x,v\rangle\bmod3,
 \quad N_0=\ker\chi_v,
 \quad N_v=N_0+\mathbb Z\frac v3.
 \label{eq:exc-vecino-indice-tres}
\end{equation}

\begin{lemma}[Structure du voisin]
\label{lem:exc-estructura-vecino}
Le caractère \(\chi_v\) est surjectif, \([N:N_0]=3\),
\(v/3\in N_0^*\), et \(N_v\) est pair et unimodulaire.
\end{lemma}

\begin{proof}
Pour toute racine d'un facteur \(A_2\), le produit avec la plus haute racine est
\(\pm1\) ou \(\pm2\) ; les coefficients de \(v\) sont congrus à un
modulo trois. Par conséquent, aucune racine n'appartient à \(N_0\) et \(\chi_v\) est
non trivial, donc surjectif. La définition du noyau donne l'indice trois et
\(\det N_0=9\). Pour \(x\in N_0\),
\(\langle x,v/3\rangle\in\mathbb Z\), tandis que
\((v/3)^2=6\). Ainsi, l'extension par \(v/3\) est entière et paire. Comme
\([N_v:N_0]=3\),
\(\det N_v=9/3^2=1\).
\end{proof}

\begin{lemma}[Minimum des nouvelles classes]
\label{lem:exc-minimo-clases}
Chaque coordonnée d'une classe locale apparaissant dans
\(N_0\pm v/3\) a pour norme minimale \(2/9\). En conséquence,
\[
 \|x\pm v/3\|^2\ge12\cdot\frac29=\frac83
 \qquad(x\in N_0).
\]
\end{lemma}

\begin{proof}
En écrivant une coordonnée de \(A_2^*\) dans la base des racines simples, les
six classes orientées pertinentes sont les translations de
\(\pm\rho/3\) par les trois classes discriminantes. Compléter les carrés dans la
forme de Gram
\(\bigl(\begin{smallmatrix}2&-1\\-1&2\end{smallmatrix}\bigr)\)
donne le minimum \(2/9\) dans les six classes. La somme orthogonale des douze
coordonnées donne l'inégalité.
\end{proof}

\begin{theorem}[Construction du réseau de Leech]
\label{thm:exc-leech}
Le voisin \(N_v\) de \eqref{eq:exc-vecino-indice-tres} est isométrique au
réseau de Leech :
\begin{equation}
 \boxed{N_v\cong\Lambda_{24}.}
 \label{eq:exc-leech}
\end{equation}
\end{theorem}

\begin{proof}
Les racines de \(N\) n'appartiennent pas à \(N_0\), et les deux nouvelles classes ne
contiennent pas de vecteurs de norme deux par le
\cref{lem:exc-minimo-clases}. Le \cref{lem:exc-estructura-vecino} prouve que
\(N_v\) est défini positif, pair, unimodulaire et de rang vingt-quatre.
L'unicité du réseau possédant ces propriétés et dépourvu de racines l'identifie à
\(\Lambda_{24}\) \citep{ConwaySloane1999}.
\end{proof}

\subsection{Isométrie d'ordre trois, orbifold et module lunaire}

La rotation de Coxeter de chaque facteur \(A_2\) induit une isométrie
\(g_c\) d'ordre trois. La stabilité du code \(C_W\) prouve que
\(g_cN=N\). Pour le repère orienté qui fixe \(v\), les translations
\begin{equation}
 y_*:=\frac{g_cv-v}{3},
 \qquad
 z_*:=\frac{g_c^{-1}v-v}{3}
 \label{eq:exc-desplazamientos-vecino}
\end{equation}
appartiennent à \(N_0\) : leurs mots discriminants sont un couple
\(c_*,-c_*\in C_W\), et
\(\langle y_*,v\rangle=\langle z_*,v\rangle=-27\). Par conséquent,
\begin{equation}
 g_cN_v=N_v.
 \label{eq:exc-preservacion-vecino}
\end{equation}
L'action n'a pas de points fixes dans \(N_v\otimes\mathbb R\), car chaque
facteur \(A_2\) possède les deux valeurs propres primitives d'ordre trois.

Soit \(\widehat g_c\) le relèvement standard à la VOA de réseau
\(V_{\Lambda_{24}}\). Ses deux espaces propres non triviaux sont de dimension
douze. L'énergie du vide du module tordu est donc
\begin{equation}
 \rho_{g_c}
 =\frac1{4\cdot3^2}
 \sum_{k=1}^{2}k(3-k)\dim\mathfrak h_{(k)}
 =\frac43.
 \label{eq:exc-peso-torcido}
\end{equation}
La congruence \(3^2\rho_{g_c}=12\equiv0\pmod3\) prouve que le relèvement
est de type zéro.

\begin{theorem}[Orbifold triadique]
\label{thm:exc-orbifold}
L'orbifold cyclique du relèvement précédent satisfait
\begin{equation}
 \operatorname{Orb}_{\mathbb Z/3}
 (V_{\Lambda_{24}},\widehat g_c)\cong V^\natural.
 \label{eq:exc-orbifold}
\end{equation}
L'automorphisme dual appartient à la classe \(3B\).
\end{theorem}

\begin{proof}
Les trois hypothèses qui ne peuvent être remplacées par une
égalité de séries ont été vérifiées : isométrie sans point fixe d'ordre trois, préservation du
voisin et type zéro. Les théorèmes d'orbifold cyclique identifient alors
l'extension holomorphe au module lunaire et classent l'automorphisme dual
\citep{ChenLamShimakura2016,AbeLamYamada2017,Carnahan2017}.
\end{proof}

La construction d'ordre deux est une réalisation parallèle. Pour le relèvement
\(\theta\) de la négation sur \(\Lambda_{24}\),
\begin{equation}
 V^\natural=(V_{\Lambda_{24}})^+
 \oplus(V_{\Lambda_{24}}^T)^+.
 \label{eq:exc-flm}
\end{equation}
Les voies d'ordre deux et trois sont identifiées par les théorèmes de VOA ; elles ne sont pas
mises en série causalement l'une dans l'autre \citep{FLM1988}.

\begin{theorem}[Publications du module lunaire]
\label{thm:exc-moonshine}
Une fois \(V^\natural\) construit, on obtient
\begin{align}
 \operatorname{Aut}(V^\natural)&\cong\mathbb M,
 \label{eq:exc-aut-monster}\\
 \operatorname{ch}_{V^\natural}(\tau)
 &=J(\tau)=j(\tau)-744
 =q^{-1}+196884q+21493760q^2+\cdots,
 \label{eq:exc-caracter-j}\\
 T_g(\tau)&=\operatorname{Tr}_{V^\natural}
 \bigl(gq^{L_0-1}\bigr),\qquad g\in\mathbb M.
 \label{eq:exc-mckay-thompson}
\end{align}
Les séries \(T_g\) sont les fonctions modulaires de genre zéro du théorème
de Moonshine et satisfont les identités de réplicabilité.
\end{theorem}

\begin{proof}
La première identification et la construction graduée proviennent de la
réalisation FLM. La coordination entre l'action graduée, les séries de
McKay--Thompson et leurs propriétés modulaires est le théorème de Moonshine de
Conway--Norton et de Borcherds
\citep{FLM1988,ConwayNorton1979,Borcherds1992}. En particulier,
\((V^\natural)_0=\mathbb C\mathbf1\) et
\((V^\natural)_1=0\) ; le coefficient \(196884=1+196883\) est la première
décomposition non triviale en dimensions de représentations du Monstre.
\end{proof}

\begin{falsifier}[Corridor exceptionnel]
\label{fals:exc-corredor}
La publication est réfutée si l'une des étapes effectives échoue :
\eqref{eq:exc-despliegue-42654}, la trace
\eqref{eq:exc-traza-fock}, le coefficient \([q]t_3=54\),
\(C_PC_P^{\mathsf T}=5I_6\), \(A_W^2=-I_6\), l'autodualité ou
l'énumérateur de \(C_W\), la parité ou l'unimodularité du recollement, l'absence
de racines du voisin, sa préservation par \(g_c\), le poids tordu
\(4/3\) ou le type zéro. Est également réfuté un récit qui utilise
\(W_{24}\) pour construire \(\Lambda_{24}\), qui confond les deux ensembles
de \(243\) par leur cardinal, ou qui fait agir le Monstre directement
sur les chiffres et les voies. L'action de \(\mathbb M\) réside dans \(V^\natural\).
\end{falsifier}
